% ═══════════════════════════════════════════════════════════════
% MT-03b: Minitest hay + muebles
% Spanisch Klasse 7 | Sequenz 3: Mi casa
% ═══════════════════════════════════════════════════════════════
% Dauer: 15 Minuten
% Punkte: 28 BE
% Inhalt: Möbelvokabeln, hay + un/una
% ═══════════════════════════════════════════════════════════════

\documentclass[11pt, a4paper]{article}

\usepackage[utf8]{inputenc}
\usepackage[T1]{fontenc}
\usepackage[ngerman]{babel}
\usepackage{helvet}
\renewcommand{\familydefault}{\sfdefault}

\usepackage[margin=2cm]{geometry}
\usepackage{enumitem}
\usepackage{tabularx}
\usepackage{booktabs}

\usepackage{xcolor}
\usepackage{tcolorbox}
\tcbuselibrary{skins}

\usepackage{fancyhdr}
\usepackage{lastpage}
\usepackage{needspace}

% FARBEN
\definecolor{spanisch-color}{HTML}{c41e3a}
\definecolor{grau-color}{HTML}{888888}
\definecolor{hellgrau-color}{HTML}{f5f5f5}

\colorlet{spanisch}{spanisch-color}
\colorlet{grau}{grau-color}
\colorlet{hellgrau}{hellgrau-color}

\newcommand{\fachfarbe}{spanisch}

\newcommand{\thema}{Minitest: hay + muebles}
\newcommand{\fach}{Spanisch}
\newcommand{\klasse}{7}
\newcommand{\maxpunkte}{28}
\newcommand{\zeit}{15 Minuten}
\newcommand{\datum}{\today}

\pagestyle{fancy}
\fancyhf{}
\renewcommand{\headrulewidth}{0.4pt}
\renewcommand{\footrulewidth}{0.4pt}

\lhead{\fach{} -- Klasse \klasse}
\chead{\textbf{MT-03b}}
\rhead{\datum}

\lfoot{\zeit}
\cfoot{}
\rfoot{Seite \thepage\ von \pageref{LastPage}}

\newcommand{\punktebox}[1]{\hfill\fbox{\small \_/#1}}
\newcommand{\luecke}[1]{\rule{#1}{0.4pt}}

\newtcolorbox{infobox}{
    colback=hellgrau,
    colframe=\fachfarbe,
    boxrule=1pt,
    arc=0pt,
    left=8pt, right=8pt, top=8pt, bottom=8pt
}

\newtcolorbox{warnbox}{
    colback=white,
    colframe=red!70!black,
    boxrule=1pt,
    arc=0pt,
    left=8pt, right=8pt, top=8pt, bottom=8pt
}

\newenvironment{aufgabe}[1]{%
    \needspace{5\baselineskip}%
    \nopagebreak[4]%
    \vspace{10pt}
    \noindent\textbf{\textcolor{\fachfarbe}{Aufgabe #1}}
    \nopagebreak[4]%
    \vspace{5pt}
    \nopagebreak[4]%
    \begin{list}{}{\leftmargin=0pt}
    \item[]
}{%
    \end{list}
}

\newcommand{\es}[1]{\textcolor{\fachfarbe}{\textbf{#1}}}

\begin{document}

% ═══════════════════════════════════════════════════════════════
% KOPFBEREICH
% ═══════════════════════════════════════════════════════════════

\begin{center}
    {\Large\bfseries\textcolor{\fachfarbe}{\thema}}
\end{center}

\vspace{5pt}

\noindent
\begin{tabularx}{\textwidth}{|X|X|X|}
\hline
\textbf{Name:} \luecke{4cm} &
\textbf{Klasse:} \klasse &
\textbf{Datum:} \luecke{2cm} \\
\hline
\end{tabularx}

\vspace{10pt}

\begin{infobox}
\textbf{Zeit:} \zeit{} | \textbf{Punkte:} \maxpunkte{} BE | \textbf{Hilfsmittel:} keine
\end{infobox}

\vspace{15pt}

% ═══════════════════════════════════════════════════════════════
% AUFGABE 1: Möbel zuordnen (8 BE)
% ═══════════════════════════════════════════════════════════════

\begin{aufgabe}{1 -- Möbel zuordnen}
Verbinde die spanischen Möbel mit der deutschen Übersetzung. \punktebox{8}
\end{aufgabe}

\begin{center}
\begin{tabular}{rcl}
\es{la mesa} & \luecke{2cm} & das Bett \\[0.8em]
\es{la silla} & \luecke{2cm} & der Tisch \\[0.8em]
\es{la cama} & \luecke{2cm} & der Schrank \\[0.8em]
\es{el armario} & \luecke{2cm} & der Stuhl \\[0.8em]
\es{el sofá} & \luecke{2cm} & die Lampe \\[0.8em]
\es{la lámpara} & \luecke{2cm} & das Sofa \\[0.8em]
\es{la estantería} & \luecke{2cm} & der Schreibtisch \\[0.8em]
\es{el escritorio} & \luecke{2cm} & das Regal \\
\end{tabular}
\end{center}

\vspace{15pt}

% ═══════════════════════════════════════════════════════════════
% AUFGABE 2: un oder una (8 BE)
% ═══════════════════════════════════════════════════════════════

\begin{aufgabe}{2 -- un oder una einsetzen}
Setze den richtigen unbestimmten Artikel ein: \es{un} oder \es{una}? \punktebox{8}
\end{aufgabe}

\begin{warnbox}
\textbf{Erinnerung:} Nach \es{hay} kommt der unbestimmte Artikel (un/una), NICHT el/la!
\end{warnbox}

\vspace{0.5em}

\noindent
a) Hay \luecke{1.5cm} mesa. \hfill
b) Hay \luecke{1.5cm} sofá. \\[0.8em]
c) Hay \luecke{1.5cm} cama. \hfill
d) Hay \luecke{1.5cm} armario. \\[0.8em]
e) Hay \luecke{1.5cm} lámpara. \hfill
f) Hay \luecke{1.5cm} silla. \\[0.8em]
g) Hay \luecke{1.5cm} escritorio. \hfill
h) Hay \luecke{1.5cm} estantería. \\

\vspace{15pt}

% ═══════════════════════════════════════════════════════════════
% AUFGABE 3: Zimmer beschreiben (12 BE)
% ═══════════════════════════════════════════════════════════════

\begin{aufgabe}{3 -- Zimmer beschreiben}
Ergänze die Sätze mit \es{hay} und dem passenden Möbelstück. Achte auf den Artikel! \punktebox{12}
\end{aufgabe}

\noindent
a) En el dormitorio \luecke{5cm}. \\
\hspace*{0.5cm}\textit{(Es gibt ein Bett.)}

\vspace{0.8em}

\noindent
b) En el salón \luecke{5cm}. \\
\hspace*{0.5cm}\textit{(Es gibt ein Sofa.)}

\vspace{0.8em}

\noindent
c) En la cocina \luecke{5cm}. \\
\hspace*{0.5cm}\textit{(Es gibt einen Tisch.)}

\vspace{0.8em}

\noindent
d) En mi habitación \luecke{5cm}. \\
\hspace*{0.5cm}\textit{(Es gibt einen Schreibtisch.)}

\vspace{0.8em}

\noindent
e) En el comedor \luecke{5cm}. \\
\hspace*{0.5cm}\textit{(Es gibt Stühle.)}

\vspace{0.8em}

\noindent
f) En el dormitorio \luecke{5cm}. \\
\hspace*{0.5cm}\textit{(Es gibt einen Schrank.)}

% ═══════════════════════════════════════════════════════════════
% PUNKTEÜBERSICHT
% ═══════════════════════════════════════════════════════════════

\vspace{25pt}
\hrule
\vspace{10pt}

\noindent
\begin{tabularx}{\textwidth}{|X|c|c|c||c|}
\hline
\textbf{Aufgabe} & 1 & 2 & 3 & \textbf{Gesamt} \\
\hline
\textbf{max. Punkte} & 8 & 8 & 12 & \textbf{28} \\
\hline
\textbf{erreicht} & & & & \\[0.8em]
\hline
\end{tabularx}

\vspace{10pt}

\begin{flushright}
\textbf{Note:} \luecke{2cm}
\end{flushright}

\end{document}
